\documentclass[12pt]{article}

\usepackage{sbc-template}
\usepackage{graphicx,url}
\usepackage[utf8]{inputenc}
\usepackage[T1]{fontenc}
\usepackage{booktabs}
\usepackage{amsmath}
\usepackage{xcolor}

\newcommand{\preencher}[1]{\textcolor{red}{\bfseries [A PREENCHER: #1]}}

\sloppy

\title{Predição de $\lambda_{max}$ de Moléculas Orgânicas via Aprendizado Profundo: Comparação de Arquiteturas e Aplicação em Triagem Computacional}

\author{Crhisângela Ferreira\inst{1}}

\address{Programa de Pós-Graduação em Ciência da Computação -- Universidade Federal do ABC (UFABC)\\
  CCM-109 -- Tópicos Especiais em Inteligência Artificial\\
  \email{\preencher{e-mail institucional ou de contato}}
}

\begin{document}

\maketitle

\begin{resumo}
  A predição do comprimento de onda de máxima absorção ($\lambda_{max}$) de moléculas
  orgânicas é relevante para o desenvolvimento de corantes, sensores e fotocatalisadores,
  mas métodos quânticos de referência (DFT) têm custo computacional elevado para triagens
  em larga escala. Este trabalho treina e compara quatro abordagens -- MLP com fingerprints
  Morgan, Gradient Boosting sobre descritores RDKit, LSTM bidirecional sobre SMILES e,
  opcionalmente, uma rede de mensagens direcionadas (D-MPNN) -- para prever $\lambda_{max}$
  a partir do SMILES do cromóforo e do solvente, usando o dataset público Deep4Chem. O
  melhor modelo é então aplicado, com quantificação de incerteza via MC-Dropout, a um
  conjunto proprietário de 5.974 moléculas geradas computacionalmente, para priorizar
  candidatos à síntese experimental. \preencher{completar o resumo com o resultado
  quantitativo principal (ex.: MAE do melhor modelo) assim que os experimentos finais
  estiverem concluídos}.
\end{resumo}

\begin{abstract}
  \preencher{traduzir o resumo para o inglês (abstract) após finalizar o "resumo" em
  português}
\end{abstract}


\section{Descrição do Problema}
\label{sec:problema}

A predição de propriedades ópticas de moléculas orgânicas é um desafio central em
quimioinformática e no desenvolvimento de novos materiais. Em particular, o comprimento
de onda de máxima absorção ($\lambda_{max}$) no espectro UV-Vis determina aplicações em
corantes, sensores e fotocatalisadores. Métodos quânticos como a Teoria do Funcional da
Densidade (DFT) permitem calcular $\lambda_{max}$ com boa precisão, mas seu custo
computacional elevado inviabiliza a triagem em larga escala sobre espaços moleculares
amplos.

Este trabalho formula o problema como uma tarefa de regressão supervisionada: dado o
SMILES de um cromóforo e o SMILES do solvente em que ele é medido, prever o valor de
$\lambda_{max}$ (em nanômetros). A escolha de incluir o solvente como entrada decorre do
fato de que $\lambda_{max}$ é sensível ao ambiente químico (efeitos solvatocrômicos), e
não apenas à estrutura do cromóforo isolado.

O projeto é dividido em duas etapas. Na \textbf{Etapa 1}, modelos são treinados e
comparados sobre um dataset público rotulado, com o objetivo de identificar a abordagem
com melhor generalização. Na \textbf{Etapa 2}, o modelo selecionado é aplicado como
ferramenta de triagem exploratória sobre um conjunto proprietário de 5.974 moléculas sem
rótulo, geradas computacionalmente por um aluno de iniciação científica do laboratório
ABC Sim da UFABC, produzindo um ranking de candidatos prioritários para síntese
experimental.

\subsection{Motivação}

A motivação prática do projeto é substituir, para fins de triagem inicial, cálculos DFT
custosos por um modelo de aprendizado de máquina treinado sobre dados experimentais, capaz
de gerar predições em escala para milhares de moléculas em segundos. Isso é particularmente
relevante para o laboratório ABC Sim, que já dispõe de um espaço de moléculas geradas
computacionalmente e precisa de um critério eficiente para priorizar quais estruturas
sintetizar e caracterizar experimentalmente.

\subsection{Definição Formal da Tarefa}

Seja $s_c$ o SMILES do cromóforo e $s_v$ o SMILES do solvente. O objetivo é aprender uma
função $f(s_c, s_v) \rightarrow \hat{y}$, em que $\hat{y}$ é a predição de $\lambda_{max}$
em nanômetros, minimizando o Erro Médio Absoluto (MAE) entre $\hat{y}$ e o valor
experimental $y$ em um conjunto de teste mantido separado do treino.


\section{Abordagem}
\label{sec:abordagem}

\subsection{Conjunto de Dados}

O treinamento e a avaliação dos modelos foram realizados sobre o \textbf{Deep4Chem}
\cite{joung2020}, base de dados pública disponível via Figshare (DOI:
10.6084/m9.figshare.12045567.v2). O Deep4Chem compila propriedades ópticas experimentais
de 7.016 cromóforos orgânicos únicos medidos em 365 solventes distintos, totalizando
20.236 registros extraídos de 1.358 artigos científicos. Cada entrada contém o SMILES do
cromóforo, o SMILES do solvente, $\lambda_{abs}$, $\lambda_{em}$, coeficiente de extinção
molar e rendimento quântico; para este trabalho, apenas o SMILES do cromóforo, o SMILES do
solvente e $\lambda_{abs}$ (renomeado \texttt{lambda\_max}) foram utilizados.

Na Etapa 2, o modelo treinado é aplicado sobre um conjunto proprietário de 5.974 moléculas,
representadas por SMILES canônicos e por colunas de identificação que rastreiam os
anéis/fragmentos de origem de cada molécula. Esse conjunto não possui rótulos de
$\lambda_{max}$ e é usado exclusivamente para triagem exploratória, nunca para treino ou
avaliação.

\subsection{Pré-processamento}

Os SMILES foram validados com o RDKit \cite{rdkit} (parsing bem-sucedido via
\texttt{Chem.MolFromSmiles}), descartando-se registros inválidos ou incompletos. Como
$\lambda_{max}$ depende do solvente, o SMILES do solvente foi mantido como feature em
todos os modelos, e não descartado no pré-processamento.

\textbf{Divisão dos dados (scaffold split).} Em vez de uma divisão aleatória entre treino
e teste, foi adotado um \emph{scaffold split} por molécula, baseado no arcabouço de
Bemis-Murcko \cite{bemis1996}: todas as ocorrências de uma mesma família estrutural
(scaffold) são mantidas inteiramente no treino ou inteiramente no teste. Essa escolha
evita que o mesmo esqueleto molecular apareça em treino e teste, o que superestimaria a
capacidade de generalização do modelo caso fosse usada uma divisão aleatória --
recomendação incorporada a partir do feedback recebido sobre a proposta inicial do
projeto.

\subsection{Modelos}

Foram implementadas e comparadas quatro abordagens, todas em PyTorch/scikit-learn,
treinadas no Google Colab (GPU gratuita):

\begin{itemize}
  \item \textbf{Baseline 1 -- MLP com fingerprints Morgan.} Os SMILES do cromóforo e do
  solvente são convertidos em vetores binários via fingerprints de Morgan
  \cite{rogers2010} (2.048 bits, raio 2, calculados com RDKit) e concatenados como entrada
  de uma rede totalmente conectada (MLP) com camadas ocultas de 512 e 256 unidades,
  ativação ReLU e \emph{dropout}.

  \item \textbf{Baseline 2 (forte) -- Gradient Boosting sobre descritores RDKit.} Um
  \texttt{GradientBoostingRegressor} \cite{friedman2001} (500 estimadores, profundidade
  máxima 4, taxa de aprendizado 0,05) é treinado sobre descritores físico-químicos
  calculados diretamente pelo RDKit (peso molecular, LogP, número de anéis aromáticos,
  entre outros), incorporado como comparação adicional sugerida no feedback da proposta.

  \item \textbf{Modelo principal -- LSTM bidirecional \emph{char-level}.} Cada caractere
  do SMILES do cromóforo é mapeado a um \emph{embedding} treinável e processado por uma
  LSTM bidirecional \cite{schuster1997}; sua representação final é concatenada ao
  fingerprint Morgan do solvente (256 bits) antes de passar por camadas densas que
  produzem a predição final.

  \item \textbf{Extensão opcional -- D-MPNN (Chemprop).} Uma rede de mensagens direcionadas
  sobre o grafo molecular \cite{yang2019}, incluída como comparação de estado da arte,
  condicionada à disponibilidade de tempo para treino (célula opcional no notebook,
  requer a biblioteca \texttt{chemprop}).
\end{itemize}

\subsection{Protocolo Experimental}

Todos os modelos foram avaliados com o mesmo scaffold split (80\% treino / 20\% teste), e
a métrica principal reportada é o Erro Médio Absoluto (MAE, em nm), complementada pela
Raiz do Erro Quadrático Médio (RMSE, em nm) -- ambas incorporadas a partir da recomendação
de reportar as duas métricas em vez de apenas o MAE. Hiperparâmetros foram selecionados
por validação cruzada sobre o conjunto de treino.

\subsection{Etapa 2 -- Aplicação e Quantificação de Incerteza}

O modelo com melhor desempenho na Etapa 1 é aplicado ao conjunto proprietário de 5.974
moléculas para gerar um ranking de $\lambda_{max}$ previsto. Como o conjunto proprietário
não possui um solvente rotulado, foi adotado um solvente de referência constante para
todas as predições -- decisão que deve ser ajustada conforme o contexto experimental real
do laboratório ABC Sim (por exemplo, o solvente mais frequente no Deep4Chem, ou o de
interesse direto do laboratório).

Como esse ranking guiará decisões de síntese experimental, foi incorporada quantificação
de incerteza via \emph{MC-Dropout} \cite{gal2016}: o modelo é mantido em modo de
treinamento (\emph{dropout} ativo) durante a inferência, e múltiplas passadas
($n=30$) sobre a mesma molécula produzem uma distribuição de predições, cujo desvio-padrão
é usado como proxy de incerteza. Candidatos de alta confiança (baixa variância entre
passadas) são priorizados no ranking final. Também está prevista uma análise de
deslocamento de distribuição (\emph{domain shift}) entre o Deep4Chem e o conjunto
proprietário, comparando estatísticas de descritores moleculares (peso molecular, LogP
etc.) entre os dois conjuntos, para avaliar até que ponto a transferência do modelo é
confiável.


\section{Resultados}
\label{sec:resultados}

\preencher{Esta seção está estruturada como template. Os experimentos completos (GBM,
Bi-LSTM, comparação final e aplicação na Etapa 2) ainda não haviam sido executados até a
elaboração deste relatório -- preencher as tabelas e figuras abaixo após rodar o notebook
por completo no Google Colab, com os datasets reais.}

\subsection{Resultado preliminar disponível}

Até o momento, apenas o treinamento do Baseline 1 (MLP com fingerprints Morgan) foi
executado, com a seguinte evolução do MAE de validação ao longo dos epochs:

\begin{table}[ht]
\centering
\caption{MAE de validação do MLP (fingerprints Morgan) por epoch -- resultado
preliminar, execução parcial.}
\label{tab:mlp-preliminar}
\begin{tabular}{cc}
\toprule
Epoch & MAE de validação (nm) \\
\midrule
20  & 36,90 \\
40  & 35,04 \\
60  & 33,91 \\
80  & 33,80 \\
100 & \preencher{valor do epoch 100} \\
\bottomrule
\end{tabular}
\end{table}

O valor de MAE de validação estabiliza em torno de 33,8--34 nm entre os epochs 80 e 100,
sugerindo convergência do baseline nessa faixa. \preencher{comentar se esse patamar é
consistente com a literatura sobre o Deep4Chem, e se há sinais de overfitting}.

\subsection{Comparação entre modelos (Etapa 1)}

\preencher{preencher com os valores de MAE e RMSE de teste, no scaffold split, para cada
modelo}

\begin{table}[ht]
\centering
\caption{Comparação de desempenho dos modelos na Etapa 1 (conjunto de teste, scaffold
split).}
\label{tab:comparacao}
\begin{tabular}{lcc}
\toprule
Modelo & MAE (nm) & RMSE (nm) \\
\midrule
MLP + fingerprints Morgan          & \preencher{valor} & \preencher{valor} \\
Gradient Boosting (descritores RDKit) & \preencher{valor} & \preencher{valor} \\
Bi-LSTM (SMILES + solvente)        & \preencher{valor} & \preencher{valor} \\
D-MPNN (Chemprop, opcional)        & \preencher{valor} & \preencher{valor} \\
\bottomrule
\end{tabular}
\end{table}

\preencher{Incluir aqui um gráfico (dispersão predito vs. real, ou barras de MAE por
modelo) gerado a partir das saídas reais do notebook -- por exemplo, com
\texttt{matplotlib}, exportado como PNG/PDF e incluído via \texttt{\textbackslash
includegraphics}.}

\subsection{Etapa 2 -- Ranking e incerteza}

\preencher{preencher com: (i) estatísticas descritivas do ranking de $\lambda_{max}$
previsto para as 5.974 moléculas; (ii) exemplos dos candidatos de maior confiança
(menor desvio-padrão via MC-Dropout); (iii) resultado da análise de deslocamento de
distribuição entre Deep4Chem e o conjunto proprietário.}


\section{Lições Aprendidas}
\label{sec:licoes}

\preencher{Esta seção deve ser escrita após a execução completa dos experimentos.
Sugestões de pontos a cobrir, com base no protocolo adotado:}

\begin{itemize}
  \item \preencher{O que funcionou: por exemplo, se o scaffold split revelou uma queda de
  desempenho em relação a um split aleatório testado como controle, evidenciando
  overfitting a famílias estruturais específicas.}
  \item \preencher{O que não funcionou ou exigiu ajuste: por exemplo, dificuldades no
  parsing de SMILES do Deep4Chem, necessidade de ajustar \texttt{RAW\_PATH}/\texttt{COL\_MAP}
  conforme o cabeçalho real do arquivo baixado do Figshare, ou instabilidade no treino do
  Bi-LSTM.}
  \item \preencher{Decisões de design que valeriam revisão: por exemplo, a escolha de um
  solvente de referência constante para a Etapa 2, dado que o conjunto proprietário não
  tem solvente rotulado -- isso é uma limitação a ser discutida.}
  \item \preencher{O que seria feito diferente: por exemplo, priorizar o D-MPNN desde o
  início em vez de deixá-lo como extensão opcional, ou aplicar validação cruzada por
  scaffold (k-fold sobre scaffolds) em vez de um único split, conforme sugerido no
  feedback recebido.}
\end{itemize}


\section{Reprodutibilidade}
\label{sec:reprodutibilidade}

O código completo do projeto (pré-processamento, treinamento dos quatro modelos,
avaliação e aplicação na Etapa 2) está disponível em:
\preencher{inserir aqui o link do repositório (GitHub/GitLab)}.

O repositório inclui o notebook utilizado para todos os experimentos, com instruções para
baixar o Deep4Chem via Figshare e para ajustar os caminhos dos arquivos de entrada.


\bibliographystyle{sbc}
\bibliography{referencias}

\end{document}
